\DocumentMetadata{
  pdfversion=2.0,
  pdfstandard=ua-2,
  lang=en-US,
  tagging=on,
  tagging-setup={math/setup=mathml-SE,tabsorder=structure}
}
\documentclass[11pt,letterpaper]{article}
\usepackage[margin=0.68in,includefoot]{geometry}
\usepackage{newcomputermodern}
\usepackage{amsmath,amscd,mathtools}
\providecommand{\square}{\mdlgwhtsquare}
\usepackage[hidelinks]{hyperref}
\hypersetup{
  pdftitle={Numerical Linear Algebra First-Year Exam, May 2022},
  pdfauthor={University of Florida Department of Mathematics},
  pdfsubject={Graduate mathematics examination},
  pdfdisplaydoctitle=true,
  bookmarksopen=true
}
\setcounter{secnumdepth}{0}
\setlength{\parindent}{0pt}
\setlength{\parskip}{0.28em}
\setlength{\emergencystretch}{3em}
\setlength{\leftmargini}{2.15em}
\setlength{\leftmarginii}{2.65em}
\setlength{\leftmarginiii}{3em}
\renewcommand{\labelenumi}{\textbf{\theenumi.}}
\pagestyle{empty}
\raggedbottom
\allowdisplaybreaks
\newenvironment{examproblems}[1]{%
  \begin{enumerate}%
  \setcounter{enumi}{#1}%
  \setlength{\topsep}{0.35em}%
  \setlength{\partopsep}{0pt}%
  \setlength{\itemsep}{0.48em}%
  \setlength{\parsep}{0.18em}%
}{\end{enumerate}}
\newenvironment{examsubparts}{%
  \begin{enumerate}%
  \setlength{\topsep}{0.18em}%
  \setlength{\partopsep}{0pt}%
  \setlength{\itemsep}{0.16em}%
  \setlength{\parsep}{0.1em}%
}{\end{enumerate}}
\newenvironment{examinstructions}{%
  \par\begingroup\itshape
}{%
  \par\endgroup\vspace{0.18em}
}
\makeatletter
\renewcommand\section{\@startsection{section}{1}{0pt}%
  {-1.25ex plus -.25ex minus -.1ex}%
  {0.55ex plus .1ex}%
  {\normalfont\large\bfseries}}
\renewcommand{\@maketitle}{%
  \newpage\null\vskip -1.2em%
  {\footnotesize\bfseries
    \href{https://www.ufl.edu/}{University of Florida}, \href{https://math.ufl.edu/}{Department of Mathematics}\par}%
  \vskip 0.18em%
  \tagstructbegin{tag=Title}%
    {\LARGE\bfseries\href{https://gma.math.ufl.edu/exams/numerical-linear-algebra-first-year/}{Numerical Linear Algebra First-Year Exam}, May 2022\par}%
  \tagstructend%
  \vskip 0.35em%
  \tagmcbegin{artifact}%
    \hrule height 1pt%
  \tagmcend%
  \vskip 0.55em%
}
\makeatother
\title{Numerical Linear Algebra First-Year Exam, May 2022}
\author{}
\date{}
\begin{document}
\maketitle
\thispagestyle{empty}
\begin{examinstructions}
Do 4 (four) problems.
\end{examinstructions}
\begin{examproblems}{0}
\item
Prove or provide a counterexample to each of the following.
\begin{examsubparts}
\item[\textnormal{(a)}]
If matrix~\(A\) is normal and triangular, then it is diagonal.
\item[\textnormal{(b)}]
Every matrix has a Schur factorization.
\end{examsubparts}
\item
This problem has two parts.
\begin{examsubparts}
\item[\textnormal{(a)}]
For \(v\in\mathbb{C}^n\), define \(f_A(v):=(v^*Av)^{1/2}\). Under what condition(s) on \(A\in\mathbb{C}^{n\times n}\) is \(f_A(\cdot)\) a norm on \(\mathbb{C}^n\)? (Full credit will only be given for a general condition or conditions. An example is not sufficient.) Prove that \(f_A(\cdot)\) is a norm on \(\mathbb{C}^n\) under the condition(s) you require.
\item[\textnormal{(b)}]
Assuming the condition(s) you require in (a), given a matrix~\(A\), determine the best constants~\(\alpha\) and~\(\beta\) to satisfy the inequality 
\[
\alpha\lVert v\rVert_2\leq f_A(v)\leq\beta\lVert v\rVert_2.
\]
\end{examsubparts}
\item
Let \(A=U\Sigma V^*\) be the singular value decomposition of \(A\in\mathbb{C}^{m\times n}\). Let \(u_j\) denote column~\(j\) of~\(U\).
\begin{examsubparts}
\item[\textnormal{(a)}]
Suppose \(\operatorname{rank}(A)=p<n<m\). Show \(\{u_1,u_2,\ldots,u_p\}\), is a basis for \(\operatorname{Col}(A)\), and \(\{u_{p+1},u_{p+2},\ldots,u_m\}\) is a basis for \(\operatorname{Null}(A^*)\).
\item[\textnormal{(b)}]
Suppose~\(A\) is full rank and \(x\neq 0\). Let \(\sigma_i\), \(i=1,\ldots,n\), be the singular values of~\(A\). Show 
\[
\sigma_1\geq\frac{\lVert Ax\rVert_2}{\lVert x\rVert_2}\geq\sigma_n>0.
\]
 If you want to use the property that \(\lVert A\rVert_2=\sigma_1\), then you must prove that also.
\end{examsubparts}
\item
Let \(\{a_1,\ldots,a_n\}\) be a linearly independent set of vectors. Consider the Gram–Schmidt and modified Gram–Schmidt algorithms for computing an orthonormal basis \(\{q_1,\ldots,q_n\}\) so that \(\operatorname{Span}\{q_1,\ldots,q_k\}=\operatorname{Span}\{a_1,\ldots,a_k\}\), for each \(k=1,\ldots,n\). Suppose in computing \(q_2\), an orthogonalization error is committed so that \(q_2^*q_1=\epsilon\).
\begin{examsubparts}
\item[\textnormal{(a)}]
Use the Gram–Schmidt algorithm to compute \(v_3\) so that \(q_3=v_3/\lVert v_3\rVert\). What is \(v_3^*q_2\)?
\item[\textnormal{(b)}]
Use the modified Gram–Schmidt algorithm to compute \(v_3\) so that \(q_3=v_3/\lVert v_3\rVert\). What is \(v_3^*q_2\)?
\end{examsubparts}
\item
Compute the Cholesky decomposition of the following matrix, or explain why it does not exist. 
\[
A=\begin{pmatrix}
1 & 1/2 & 2 & 3\\
1/2 & 5/16 & 3/2 & 5/2\\
2 & 3/2 & 17 & 17\\
3 & 5/2 & 17 & 31
\end{pmatrix}.
\]
\end{examproblems}
\end{document}
